\documentclass[uplatex,a4paper,11pt]{jsarticle}

% ===== パッケージ =====
\usepackage[utf8]{inputenc}
\usepackage[dvipdfmx]{graphicx}
\usepackage{amsmath, amssymb, amsthm}
\usepackage{geometry}
\usepackage{tikz}
\usetikzlibrary{automata, positioning}
\usepackage{hyperref}
\usepackage{here}
\usepackage{bookmark}

\geometry{margin=25mm}
\renewcommand{\baselinestretch}{1.3}

% ===== 定理環境 =====
\newtheorem{theorem}{定理}
\newtheorem{lemma}{補題}
\newtheorem{definition}{定義}
\newtheorem{proposition}{命題}

% ===== コマンド =====
\newcommand{\N}{\mathbb{N}}
\newcommand{\Z}{\mathbb{Z}}

% ===== タイトル =====
\title{Homework 1a}
\author{学籍番号: 1270374 \\ 名前: 松本 奈生}
\date{\today}

\begin{document}

\maketitle

\section*{Question 1}
Let A and B be sets of non-negative integers defined as follows:

$
      A = \{2^k | k \in N_0 \quad and \quad k \le 7\}
$,

$
      B = \{n | n \in N_0 \quad and \quad n \le 255\}
$,

where $N_0 = N \cup \{0\}$. 
Let $f$ be a function defined as follows:

$f:P(A) \to B \quad and \quad f(X) = \sum_{a\in X} a$

where P(A) is the power set of A and $\sum_{a\in X} a$ is the 
summation of all element $a$ in $X$.

\subsection*{問題}
\begin{enumerate}
      \item Q : List all the elements of A. \\
            A : $A = \{1,2,4,8,16,32,64,128\}$
      \item Q : Answer the number of elements of P(A). \\
            A : $n(P(A)) = 2^8 = 256$
      \item Q : Answer the value of $f(\{1, 4\})$. \\
            A : $f(\{1,4\}) = 5$
      \item Q : Answer all X such that $f(X) = 10$. \\
            A : $X = \{2,8\}$
\end{enumerate}

\section*{Question 2}
Let $\Sigma = \{a,b,c\}$ and consider the following languages over $\Sigma$:


$L_1 = \{w | w \ \text{is a string over}  \ \Sigma  \  \text{and} \ |w| \le 5\},$

$L_2 = \{w | w \in L_1 \ \text{and} \ w \ \text{contains the substring} \ a b a \},$

$L_3 = \{w | w \in L_1 \ \text{and} \ w \ \text{is a palindrome}\},$

$L_4 = \{xy | x \in L_2, y \in L_3 \ \text{and} \ |y| = 3 \},$

\subsection*{問題}
\begin{enumerate}
      \item Q : Answer the number of elements of $L_1$. \\
            A : $n(L_1) = \sum_{i = 0}^{5}{3^i} = 364$
      \item Q : Answer the number of elements of $L_2$. \\
            A : $n(L_2) = 1 + 6 + 26 = 33$
      \item Q : List the first 10 elements of $L_3$ in the lexicographic order. \\
            A : $\{\epsilon, \{a\}, \{b\}, \{c\}, \{aa\}, \{bb\}, \{cc\}, \{aaa\}, \{aba\}, \{aca\}\}$
      \item Q : List the last 5 elements of $L_3$ in the lexicographic order. \\
            A : $\{\{cbbbc\}, \{cbcbc\}, \{ccacc\}, \{ccbcc\}, \{ccccc\}\}$
      \item Q : Answer the number of elements of $L_2 \cap L_3$. \\
            A : $n(L_2 \cap L_3) = 5$
      \item Q : Answer the number of elements of $L_4$. \\
            A : $n(x) = 33, n(y) = 9$より, \\ \quad \; $n(L_4) = 33 \times 9 = 297$
\end{enumerate}

\section*{Question 3}
Let $G_1 = (V_1, E_1)$ be the undirected graph shown below.

\begin{figure}[H]
      \centering
      \includegraphics[width=0.7\linewidth]{matching.png.pdf}
      \caption{matching}
      \label{fig:sample}
\end{figure}

\subsection*{問題}
\begin{enumerate}
      \item Q : Write a formal description of $G_1$. \\
            A : $G_1 = (\{1,2,3,4,5,6,7,8\}, \{(1,6),(1,8),(2,7),(2,8),(3,5),(3,6),(3,7),(4,5),(4,8)\})$
      \item Q : A spanning subgraph is a subgraph that retains every vertex of the original graph. A perfect matching is a spanning subgraph in which the degree of every vertex is exactly one. Answer all the perfect matching of $G_1$
                  (either in graphical drawing or in formal description).  \\
            A : $(\{1,2,3,4,5,6,7,8\}, \{(1,6),(2,7),(3,5),(4,8)\})$ \\
                  \quad \ \ $(\{1,2,3,4,5,6,7,8\},\{(1,8),(2,7),(3,6),(4,5)\})$
      \item Q : Answer whether of not the cnf-formula \\
                  $\phi = (x_{35} \lor x_{36} \lor x_{37}) \land (\neg x_{35} \lor \neg x_{36}) \land (\neg x_{36} \lor \neg x_{37}) \land (\neg x_{35} \lor \neg x_{37})$ \\
                  is satisifiable. If the answer is yes, then also answer an example of the truth assignments that make $\phi = 1$. \\
            A : $(x_{35}, x_{36}, x_{37}) = (1,0,0)$で充足可能.
\end{enumerate}

\section*{Question 4}
The lexicographic order defined in this textbook is different from the usual dictionary order. For example, "$appropriate$"
precedes "$back$" 
in dictionaries, while $back < appropriate$
in the lexicographic order. Why does the textbook define the order of strings as so? What is the purpose or the merit of the lexicographic order?

\subsection*{解答}
The reason for sorting strings by size in ascending order is that, naturally, strings with smaller sizes have fewer possible combinations. Therefore, when searching through a string sorted in lexicographical order, starting with the smaller strings offers the advantage of more efficient searching.
For example, let k be the size of the character we want to find.
When searching a string of n characters sorted in lexicographical order,
the time complexity can be expressed as $\sum_{i=0}^{k}{O(n^i)}$,
and since this results in a constant time complexity for any value of k,

文字列のサイズを小さい順に並べる理由は、サイズの小さい文字列の方が当然、組み合わせの総数が少ないため、辞書式順序でソートされた文字列を探索するときに小さいサイズの文字列から探索することで効率的に探索することができる利点があるから。
例えば、探索したい文字のサイズをkとして、
使用される文字数nの文字列を辞書式順序でソートして探索する時、
$\sum_{i=0}^{k}{O(n^i)}$と考えることができ、
任意のkに対して、一定のオーダーを取ることができるから。

\end{document}
